\section{Floor and Ceil in a Sorted Array}
\label{sec:bs-floor-ceil}

\problemstatement{We are given a sorted array $\textit{arr}$ of $n$
elements and a value $X$. We must find the \textbf{floor} of $X$ -- the
largest element in $\textit{arr}$ that is $\le X$ -- and the
\textbf{ceil} of $X$ -- the smallest element in $\textit{arr}$ that is
$\ge X$. If no floor (resp.\ ceil) exists, report $-1$.}

\begin{patternbox}
\emph{``Floor''} and \emph{``ceil''} are simply ``largest $\le X$'' and
``smallest $\ge X$'' -- and ``smallest $\ge X$'' is, value-for-value, the
\emph{value} sitting at the lower-bound index from
Section~\ref{sec:bs-lower-bound}. ``Largest $\le X$'' is its mirror image,
and is found with the symmetric boundary-search template, narrowing toward
the right instead of the left. Whenever a problem uses the words
\emph{``floor''} or \emph{``ceil''} (or, equivalently, \emph{``closest
element not exceeding''} / \emph{``closest element not less than''}),
translate it immediately into a boundary search.
\end{patternbox}

\subsection*{Understanding the problem carefully}

\textbf{Ceil} of $X$: define $P(i) = (\textit{arr}[i] \ge X)$. As in lower
bound, find the smallest index where $P$ holds; the ceil is the value at
that index (or $-1$ if no such index exists, i.e.\ $X$ exceeds every
element).

\textbf{Floor} of $X$: define $Q(i) = (\textit{arr}[i] \le X)$. This
predicate is \texttt{true} for a \emph{prefix} of the array and
\texttt{false} for the remaining suffix -- the mirror image of $P$. We want
the \emph{largest} index where $Q$ holds, so instead of narrowing left
after a success (as in lower/upper bound), we narrow \emph{right} after a
success, since we are now searching for the rightmost true position rather
than the leftmost one.

\begin{ideabox}[Candidate Space and Predicate]
For ceil: candidate space is indices $0,\dots,n-1$ (or ``not found''),
predicate $P(i) = (\textit{arr}[i] \ge X)$, want smallest true index.
For floor: same candidate space, predicate $Q(i) = (\textit{arr}[i] \le X)$,
want \emph{largest} true index.
\end{ideabox}

\subsection*{Why floor needs the mirrored template}

In every boundary search so far (lower bound, upper bound), we wanted the
\emph{leftmost} index satisfying a predicate that becomes true only after
some crossover point, so a successful check let us narrow the window to the
\emph{left}, hunting for an even earlier success. Floor flips this: the
predicate $\textit{arr}[i] \le X$ is true \emph{before} the crossover and
false after, so we want the \emph{rightmost} successful index, meaning a
successful check should narrow the window to the \emph{right}, hunting for
an even later success that is still within the true region. This mirrored
narrowing direction is the only structural change; everything else about
the binary search skeleton is identical.

\subsection*{Algorithm}

\textbf{Ceil:}
\begin{enumerate}
  \item $\textit{lo} \gets 0$, $\textit{hi} \gets n-1$,
        $\textit{ceil} \gets -1$.
  \item While $\textit{lo} \le \textit{hi}$: $\textit{mid} \gets
        \textit{lo}+(\textit{hi}-\textit{lo})/2$. If
        $\textit{arr}[\textit{mid}] \ge X$: $\textit{ceil} \gets
        \textit{arr}[\textit{mid}]$, $\textit{hi} \gets \textit{mid}-1$.
        Else: $\textit{lo} \gets \textit{mid}+1$.
  \item Return \textit{ceil}.
\end{enumerate}

\textbf{Floor:}
\begin{enumerate}
  \item $\textit{lo} \gets 0$, $\textit{hi} \gets n-1$,
        $\textit{floor} \gets -1$.
  \item While $\textit{lo} \le \textit{hi}$: $\textit{mid} \gets
        \textit{lo}+(\textit{hi}-\textit{lo})/2$. If
        $\textit{arr}[\textit{mid}] \le X$: $\textit{floor} \gets
        \textit{arr}[\textit{mid}]$, $\textit{lo} \gets \textit{mid}+1$.
        Else: $\textit{hi} \gets \textit{mid}-1$.
  \item Return \textit{floor}.
\end{enumerate}

\begin{algorithm}[H]
\caption{Floor and Ceil}
\begin{algorithmic}[1]
\Procedure{FloorCeil}{$\textit{arr}, X$}
  \State $\textit{lo} \gets 0$, $\textit{hi} \gets n-1$
  \State $\textit{floor} \gets -1$, $\textit{ceil} \gets -1$
  \While{$\textit{lo} \le \textit{hi}$}
    \State $\textit{mid} \gets \textit{lo} + (\textit{hi}-\textit{lo})/2$
    \If{$\textit{arr}[\textit{mid}] \le X$}
      \State $\textit{floor} \gets \textit{arr}[\textit{mid}]$, $\textit{lo} \gets \textit{mid}+1$
    \Else
      \State $\textit{hi} \gets \textit{mid}-1$
    \EndIf
  \EndWhile
  \State (reset $\textit{lo}, \textit{hi}$ and repeat symmetrically for ceil)
  \State \Return $(\textit{floor}, \textit{ceil})$
\EndProcedure
\end{algorithmic}
\end{algorithm}

\subsection*{Dry run}

\dryrun{Take $\textit{arr} = [3, 4, 4, 7, 8, 10]$, $X = 5$.}

\textbf{Ceil} (smallest $\ge 5$):
\begin{itemize}
  \item $\textit{lo}=0,\textit{hi}=5$: $\textit{mid}=2$,
        $\textit{arr}[2]=4 \ge 5$? No. $\textit{lo}=3$.
  \item $\textit{lo}=3,\textit{hi}=5$: $\textit{mid}=4$,
        $\textit{arr}[4]=8 \ge 5$? Yes. $\textit{ceil}=8$, $\textit{hi}=3$.
  \item $\textit{lo}=3,\textit{hi}=3$: $\textit{mid}=3$,
        $\textit{arr}[3]=7 \ge 5$? Yes. $\textit{ceil}=7$, $\textit{hi}=2$.
  \item $\textit{lo}=3 > \textit{hi}=2$: loop ends. $\textit{ceil}=7$.
\end{itemize}

\textbf{Floor} (largest $\le 5$):
\begin{itemize}
  \item $\textit{lo}=0,\textit{hi}=5$: $\textit{mid}=2$,
        $\textit{arr}[2]=4 \le 5$? Yes. $\textit{floor}=4$, $\textit{lo}=3$.
  \item $\textit{lo}=3,\textit{hi}=5$: $\textit{mid}=4$,
        $\textit{arr}[4]=8 \le 5$? No. $\textit{hi}=3$.
  \item $\textit{lo}=3,\textit{hi}=3$: $\textit{mid}=3$,
        $\textit{arr}[3]=7 \le 5$? No. $\textit{hi}=2$.
  \item $\textit{lo}=3 > \textit{hi}=2$: loop ends. $\textit{floor}=4$.
\end{itemize}

So floor$(5) = 4$ and ceil$(5) = 7$.

\subsection*{C++ implementation}

\begin{lstlisting}
#include <bits/stdc++.h>
using namespace std;

pair<int,int> floorCeil(vector<int>& arr, int X) {
    int n = arr.size();
    int floorVal = -1, ceilVal = -1;

    // Floor: largest element <= X
    int lo = 0, hi = n - 1;
    while (lo <= hi) {
        int mid = lo + (hi - lo) / 2;
        if (arr[mid] <= X) {
            floorVal = arr[mid];
            lo = mid + 1;
        } else {
            hi = mid - 1;
        }
    }

    // Ceil: smallest element >= X
    lo = 0; hi = n - 1;
    while (lo <= hi) {
        int mid = lo + (hi - lo) / 2;
        if (arr[mid] >= X) {
            ceilVal = arr[mid];
            hi = mid - 1;
        } else {
            lo = mid + 1;
        }
    }

    return {floorVal, ceilVal};
}

int main() {
    vector<int> arr = {3, 4, 4, 7, 8, 10};
    auto [floorVal, ceilVal] = floorCeil(arr, 5);
    cout << "Floor: " << floorVal << ", Ceil: " << ceilVal << endl;
    return 0;
}
\end{lstlisting}

\begin{complexitybox}
\textbf{Time Complexity.} Two independent binary searches, each
$O(\log n)$, giving an overall time complexity of
\[
  O(\log n).
\]
\textbf{Space Complexity.} $O(1)$ auxiliary space.
\end{complexitybox}

\begin{takeawaybox}
Floor and ceil are the value-returning siblings of lower and upper bound.
Whichever direction the predicate's ``true'' region lies in (a prefix or a
suffix of the sorted array) tells you which direction to narrow the search
window after a successful check: narrow toward the \emph{far} end of the
true region if you want its near boundary, or keep the running answer
updated while narrowing toward the \emph{near} end if you want the boundary
closest to where you started.
\end{takeawaybox}
